\subsection{PRODUCTS OF MODULES}

Let \((\mathbf{E}_i)_{i\in I}\) be a family of modules over the same ring A. It is immediately verified that the product of the module structures on the \(\mathbf{E}_i\) (I, §4, no. 8) is an A-module structure on the product set \(\mathbf{E} = \prod_{i\in I}\mathbf{E}_i\) . With this structure the set E is called the product module of the modules \(\mathbf{E}_i\) ; if \(x = (x_i)\) , \(y = (y_i)\) are two elements of E, then

\[
\left\{ \begin{array}{l l} x + y = (x _ {i} + y _ {i}) \\ \lambda . x = (\lambda x _ {i}) & \text { for   all } \lambda \in \mathbf {A}. \end{array} \right.\tag{17}
\]

Formulae (17) express the fact that the projections \(pr_{i}:E\to E_{i}\) are linear mappings; these mappings are obviously surjective.

Recall that if the indexing set I is empty, the product set \(\prod_{i\in I}E_{i}\) then consists of a single element; the product module structure on this set is then that under which this unique element is 0.

PROPOSITION 4. Let \(\mathbf{E} = \prod_{i\in I}\mathbf{E}_i\) be the product of a family of A-modules \((\mathbf{E}_i)_{i\in I}\) . For every A-module \(\mathbf{F}\) and every family of linear mappings \(f_i:\mathbf{F}\to \mathbf{E}_i\) there exists one and only one mapping \(f\) of \(\mathbf{F}\) into \(\mathbf{E}\) such that \(\mathbf{pr}_i\circ f = f_i\) for all \(i\in I\) and this mapping is linear.

This follows directly from the definitions.

Product of modules is ``associative'': if \((J_{\lambda})_{\lambda \in L}\) is a partition of I, the canonical mapping

\[
\prod_ {i \in I} E _ {i} \rightarrow \prod_ {\lambda \in L} \left(\prod_ {i \in J _ {\lambda}} E _ {i}\right)
\]

is an isomorphism.

PROPOSITION 5. (i) Let \((\mathbf{E}_t)_{i \in I}\) , \((\mathbf{F}_t)_{i \in I}\) be two families of A-modules with the same indexing set I; for every family of linear mappings \(f_i: \mathbf{E}_t \to \mathbf{F}_t\) ( \(i \in I\) ), the mapping \(f: (x_i) \to (f_i(x_i))\) of \(\prod_t \mathbf{E}_t\) into \(\prod_t \mathbf{F}_t\) (sometimes denoted by \(\prod_t f_i\) ) is linear.

\begin{enumerate}
\def\labelenumi{(\roman{enumi})}
\setcounter{enumi}{1}
\tightlist
\item
  Let \((\mathbf{G}_{\iota})_{\iota \in \mathbf{I}}\) be a third family of A-modules with I as indexing set and, for all \(t \in I, \text{ let } g_t: F_t \to G_t \text{ be a linear mapping; let } g = \prod_t g_t. \text{ For each of the sequences } E_t \xrightarrow{f_t} F_t \xrightarrow{g_t} G_t \text{ to be exact, it is necessary and sufficient that the sequence}\)
\end{enumerate}

\[
\prod_ {i} \mathrm{E} _ {i} \xrightarrow {f} \prod_ {i} \mathrm{F} _ {i} \xrightarrow {g} \prod_ {i} \mathrm{G} _ {i}
\]

be exact.

Assertion (i) follows immediately from the definitions. On the other hand, to say that \(y = (y_{i})\) belongs to \(\operatorname{Ker}(g)\) means that \(g_{i}(y_{i}) = 0\) for all \(i \in I\) and hence that \(y_{i} \in \operatorname{Ker}(g_{i})\) for all \(i \in I\) ; similarly, to say that \(y\) belongs to \(\operatorname{Im}(f)\) means that there exists \(x = (x_{i}) \in \prod_{i} E_{i}\) such that \(y = f(x)\) , which is equivalent to saying that \(y_{i} = f_{i}(x_{i})\) for all \(i \in I\) or also that \(y_{i} \in \operatorname{Im}(f_{i})\) for all \(i \in I\) ; whence (ii).

COROLLARY. Under the conditions of Proposition 5, (i),

\[
\operatorname{Ker} (f) = \prod_ {i \in I} \operatorname{Ker} (f _ {i}), \quad \operatorname{Im} (f) = \prod_ {i \in I} \operatorname{Im} (f _ {i})\tag{18}
\]

and there are canonical isomorphisms

\[
\operatorname{Coim} (f) \cong \prod_ {i \in I} \operatorname{Coim} (f _ {i}), \quad \operatorname{Coker} (f) = \prod_ {i \in I} \operatorname{Coker} (f _ {i})
\]

obtained by respectively associating with the class of an element \(x = (x_{i})\) of \(\prod_{i} E_{i}\) , mod \(\operatorname{Ker}(f)\) , (resp. with the class of an element \(y = (y_{i})\) of \(\prod_{i} F_{i}\) , mod. \(\operatorname{Im}(f)\) ) the family of classes of the \(x_{i}\) mod. \(\operatorname{Ker}(f_{i})\) (resp. the family of classes of the \(y_{i}\) mod. \(\operatorname{Im}(f)\) ).

In particular, for \(f\) to be injective (resp. surjective, bijective, zero) it is necessary and sufficient that, for all \(\iota \in I\) , \(f_{\iota}\) be injective (resp. surjective, bijective, zero).

If, for all \(\iota \in I\) , we consider a submodule \(F_{\iota}\) of \(E_{\iota}\) , the module \(\prod_{\iota \in I} F_{\iota}\) is a submodule of \(\prod_{\iota \in I} E_{\iota}\) and by virtue of the Corollary to Proposition 5, there is a canonical isomorphism

\[
\prod_ {i \in I} \left(E _ {i} / F _ {i}\right)\cong \left(\prod_ {i \in I} E _ {i}\right) / \left(\prod_ {i \in I} F _ {i}\right).\tag{19}
\]

An important example of a product of modules is that where all the factor modules are equal to the same module F; their product \(F^{I}\) is then just the set of mappings of I into F. The diagonal mapping \(F \rightarrow F^{I}\) mapping \(x \in F\) to the constant function equal to x on I is linear. If \((\mathbf{E}_{t})_{t \in I}\) is a family of A-modules and, for all \(\iota \in I\) , \(f_{\iota}: F \to E_{\iota}\) is a linear mapping, then the linear mapping \(x \mapsto (f_{\iota}(x))\) of \(F\) into \(\prod_{\iota \in I} E_{\iota}\) is the composition of the mapping \((x_{\iota}) \mapsto (f_{\iota}(x_{\iota}))\) of \(F^{\mathrm{I}}\) into \(\prod_{\iota \in I} E_{\iota}\) and the diagonal mapping \(F \to F^{\mathrm{I}}\) .

